\section{Conclusion}
A scheduled-date curve model has to answer three structural questions: what happens when information arrives, how past shocks are stored, and how meeting effects interact with the rest of the curve. The bond-price drift and announcement conditions link each answer to the chosen maturity shapes.

For announcements, a prescribed correction across maturities determines a conditional exponential moment of the level change. A linear correction leads to a Gaussian law, with the appropriate change of measure on later intervals. For continuous dynamics, a finite recurrence of the meeting loadings permits past shocks to be stored in a state whose size does not grow with the calendar. Without such a recurrence, an exact state of fixed size is impossible in the constant-scale setting studied here. On a finite calendar, a recurrent approximation can nevertheless be accurate, and the loading error supplies explicit bounds for forward rates and bond prices.

For a meeting-related curve added to a smooth component, covariance matters through the drift. The change in volatility at each meeting must be uncorrelated with the nonconstant smooth coordinates, and the remaining smooth drift restriction must use the combined level noise. These conditions give a direct way to check a proposed combination and to construct its front-end drifts. The Svensson calculation makes the permitted extra level correlation explicit. Allowing positive decay parameters to move does not remove the standalone block restrictions in an uncorrelated splice: its entire residual must still vanish.

The three-exponent construction supplies a positive stochastic-volatility example: variance is recovered from three forward rates, and discounted bonds are true martingales. It uses fixed exponents, one maturity breakpoint, and no diffusion after that date. This complements the restrictions without establishing a general classification of curve-dependent variance laws.

The examples show why exact and numerical questions should be examined together. A discrete hike/no-hike announcement is exactly excluded by the piecewise-affine jump family, but its required correction can be almost identical to the Gaussian one at realistic rate scales. A loading sequence can fail every exact finite recurrence yet admit an accurate low-order approximation on a given calendar. These distinctions make the structural results useful both for selecting a model family and for deciding how much complexity an implementation needs.

Three questions remain. The smallest autonomous state dimension under our calendar convention is not determined. The finite-calendar announcement construction does not establish a uniform-state realization theory with scheduled jumps. Finally, the numerical examples are synthetic: market calibration and assessment against observed prices remain to be carried out.
